% 4.1 — Environnement
\section{Environnement de réalisation}
\label{sec:environnement}

Le tableau~\ref{tab:env} résume l'environnement dans lequel \demo{} a été développé et testé.

\begin{table}[htbp]
\centering
\caption{Environnement et outils}
\label{tab:env}
\small
\begin{tabularx}{\linewidth}{@{}L{3.6cm}Y@{}}
\toprule
\textbf{Élément} & \textbf{Choix} \\
\midrule
Langage & Python $\geq$ 3.9, bibliothèque standard uniquement (plus \cmd{tomli} avant 3.11) ;
  aucune dépendance réseau, aucune télémétrie \\
Dépendances optionnelles & \cmd{winfspy} (montage Windows, avec le pilote WinFsp),
  \cmd{mitmproxy} (couche réseau) \\
Systèmes & Linux (x86-64) ; Windows avec PowerShell et Git Bash \\
Agents intégrés & Claude Code (Bash, PowerShell, édition de fichiers), Codex, Cursor ;
  piège bash pour le terminal \\
Tests & \cmd{pytest} ; les mêmes tests sur Linux et Windows, avec des attentes exactes par
  système quand le comportement diffère \\
Gestion de versions & git ; un commit par partie, dont le message est le bilan \\
Distribution & paquet Python (\cmd{pyproject.toml}), scripts \cmd{install.sh} et
  \cmd{install.ps1}, licence MIT \\
\bottomrule
\end{tabularx}
\end{table}

\begin{pourquoi}[Pourquoi la bibliothèque standard seulement ?]
\demo{} s'exécute à chaque commande de l'agent et se place sur le chemin de toutes ses
actions ; chaque dépendance est du code de plus qui tourne avec les droits du développeur
et une source de panne au démarrage. Le cœur n'en a aucune. Même \cmd{winfspy}, réellement
nécessaire au montage Windows, est une dépendance \emph{optionnelle} : son installation
peut échouer sans WinFsp, et cet échec ne doit pas emporter le reste de l'outil
(\cmd{pyproject.toml}).
\end{pourquoi}

\subsection{Organisation du code}
\begin{terminal}[arborescence (extrait)]
src/demo_cli/
  cli.py              # commandes : check, undo, log, verify, doctor, sh, mount...
  guard.py            # l'orchestrateur : Guard.evaluate, evaluate_file_edit
  decide.py           # le décideur unique (fonction pure)
  policy.py           # la politique : surfaces, canaux ask/deny, indices
  model.py            # l'analyse POSIX unique
  ps_model.py         # l'analyse PowerShell (C3)
  classify.py         # l'ancien chemin (règles textuelles), en retrait
  recovery.py         # captures, groupes, restauration
  checkpoint.py       # points de contrôle (optionnels)
  receipts.py         # registre chaîné SHA-256
  syscall_guard.py    # filet Linux (ptrace)
  fsguard.py          # filet Windows : le jugement
  fsmount.py          # filet Windows : WinFsp
  protect.py          # déplacement et verrouillage du projet (Windows)
  schedule.py         # tâche d'ouverture de session (W-SEC)
  hooks/              # claude_code.py, codex.py, cursor.py
  engine/             # simulate, words, invocation, effects, verdict,
                      # rules*.py (12 familles), ps_*.py, *_ops.py, *_data.py
tests/                # 138 fichiers ; fixtures/ : lignes du portail par partie
\end{terminal}
